提案手法は視覚情報のみでは把握困難な環境変動に対しても，接触時のセンサ応答に基づき把持動作を適応的に調整可能であることを示した．特筆すべきは，用いた触覚情報が抵抗値変化という低次元かつ物理的な意味が一意に定まらないデータであったにもかかわらず，模倣学習によって視覚情報や動作文脈と適切に統合され，タスク遂行に必要な接触検知や動作トリガーとして機能した点である．この結果は，高価な力覚センサによる明示的な状態推定を行わずとも，単純なセンサと学習モデルの組み合わせにより，実環境の不確実性に対してロバストな適応が可能であることを実証するものである．\par

一方で，獲得された方策の汎化能力は教示データの運動学的範囲に依存しており，教示時の可動域を大きく逸脱するような幾何学的制約に対しては，触覚情報の活用のみでは補正しきれない限界も確認された．これは，学習ベースの手法における分布外データへの対応という課題を改めて浮き彫りにするものである．